% !TEX root = ../thesis.tex
% ==============================================================================
%  CHAPTER 1: INTRODUCTION (فصل اول: مقدمه) (chapters/chapter1_introduction.tex)
% ==============================================================================
%  This chapter introduces the research problem, problem formulation,
%  significance of the study, research objectives, novel contributions,
%  and the overall roadmap of the thesis.
%
%  TYPESETTING GUIDELINES:
%  - Equation references: Use `\eqref{eq:...}` or `\ref{sec:...}`, both are adapted
%  - Chapter/Section/Figure/Table references: Use `\ref{chap:...}`, `\ref{sec:...}`,
%                                                 `\ref{fig:...}`, `\ref{tab:...}`
%  - Citations: Always prepend with a non-breaking space, e.g., `متن~\cite{ref}`
%  - English Footnotes: Use `\enfootnote{English Term}` for standard LTR footnotes
% ==============================================================================

\chapter{مقدمه}\label{chap:chapter1_introduction}

در این فصل به معرفی مسئله، اهمیت موضوع، پیشینه و چالش‌های موجود، رویکرد پیشنهادی، اهداف پژوهش و ساختار کلی پایان‌نامه پرداخته می‌شود.


\section{تعریف مسئله}\label{sec:1-1}

با وجود پیشرفت‌های چشمگیر و پذیرش گسترده‌ی مدل‌های زبانی بزرگ\enfootnote{large language models} (\lr{LLM}) در کاربردهای گوناگون، این سیستم‌ها همچنان ماهیتی به‌شدت پیچیده و غیرشفاف دارند. قدرتمندترین مدل‌های امروزی غالباً به‌صورت سیستم‌های تجاری و بسته توسعه می‌یابند که تنها از طریق رابط‌های برنامه‌نویسی (\lr{API}) قابل‌دسترسی هستند. در این سیستم‌ها که به‌عنوان توابع مولد ناشناخته ($\mathbf{y} = f(\mathbf{x}; \theta)$) عمل می‌کنند، معماری داخلی، ماتریس وزن‌ها ($\theta$) و گرادیان‌ها در پسِ رابط برنامه‌نویسی پنهان شده‌اند. این سطح از انسداد سبب می‌شود تا سازوکار تولید یک خروجی معین برای تحلیل‌گر ناشناخته باقی بماند.

به بیان ریاضی، فرض کنید $\mathcal{V}$ یک الفبای متناهی از توکن‌های گسسته‌ی واژگانی باشد ($|\mathcal{V}| < \infty$). مجموعه‌ی تمامی دنباله‌های متناهی از توکن‌ها بر روی این الفبا با بستار کلینی\enfootnote{Kleene closure} $\mathcal{V}^* = \bigcup_{t=0}^\infty \mathcal{V}^t$ تعریف می‌شود. یک پرامپت متنی ورودی، دنباله‌ای از $T_{\text{\lr{in}}}$ توکن به فرم $\mathbf{x} = (u_1, u_2, \dots, u_{T_{\text{\lr{in}}}}) \in \mathcal{V}^{T_{\text{\lr{in}}}} \subset \mathcal{V}^*$ است. یک مدل زبانی بزرگ خودهمبسته\enfootnote{autoregressive} که با بردار پارامترهای مجهول $\theta \in \Theta \subset \mathbb{R}^P$ (که در آن $P$ نمایانگر ابعاد فضای پارامترها و در مرتبه‌ی صدها میلیارد است) پیکربندی شده است، یک توزیع احتمال شرطی بر روی فضای دنباله‌های خروجی $\mathcal{V}^*$ تعریف می‌کند. با بهره‌گیری از قاعده‌ی ضرب احتمالات، توزیع احتمال شرطی تولید توالی پاسخ $\mathbf{y} = (v_1, v_2, \dots, v_{T_{\text{\lr{out}}}}) \in \mathcal{V}^{T_{\text{\lr{out}}}}$ به‌صورت زیر تجزیه می‌شود~\cite{vaswani2017attention, brown2020language}:
\begin{equation}\label{eq:chap1_autoregressive_prob}
    P_\theta(\mathbf{y} \mid \mathbf{x}) = \prod_{t=1}^{T_{\text{\lr{out}}}} P_\theta(v_t \mid \mathbf{x}, v_{<t})
\end{equation}
که در آن $v_{<t} = (v_1, \dots, v_{t-1})$ پیشوند توکن‌های تولیدشده تا پیش از گام $t$ است و $v_{<1} = \emptyset$ در نظر گرفته می‌شود. تحت یک راهبرد رمزگشایی\enfootnote{decoding} معین نظیر رمزگشایی حریصانه\enfootnote{greedy decoding}، نمونه‌برداری هسته‌ای\enfootnote{nucleus sampling}~\cite{holtzman2020curious} یا جستجوی پرتو\enfootnote{beam search}، مدل زبانی به‌عنوان یک نگاشت عمل می‌کند:
\begin{equation}\label{eq:chap1_generative_map}
    f_\theta: \mathcal{V}^* \to \mathcal{V}^*
\end{equation}
در رژیم دسترسی جعبه‌سیاه\enfootnote{black-box} (یا مدل پرس‌وجوی مرتبه‌صفر\enfootnote{zeroth-order query model}~\cite{nesterov2017random})، تحلیل‌گر هیچ‌گونه دسترسی به فضای پارامترها $\Theta$، گراف محاسباتی $\mathcal{G}$، لایه‌های پنهان یا فعال‌سازی‌های میانی $\mathbf{h}^{(l)} \in \mathbb{R}^{T \times d_{\text{\lr{model}}}}$ و عملگرهای مشتق‌پذیری و گرادیان‌ها ($\nabla_\theta$ و $\nabla_{\mathbf{x}}$) ندارد. تنها عمل مجاز برای مفسر، استعلام پرامپت $\mathbf{x}$ از اوراکل\enfootnote{oracle} سیستم و مشاهده‌ی متن گسسته‌ی محقق‌شده $\mathbf{y} = f_\theta(\mathbf{x})$ یا تحقق متغیر تصادفی $\mathbf{y} \sim P_\theta(\cdot \mid \mathbf{x})$ است، بدون آنکه حتی بردار احتمالات خروجی (لوجیت‌ها\enfootnote{logits}) در دسترس قرار گیرد.

فراتر از محدودیت‌های دسترسی، مسئله‌ی تفسیرپذیری\enfootnote{interpretability} در مدل‌های مولد مدرن از نظر ریاضی و مفهومی ذاتاً بدتعریف است. برخلاف مدل‌های متمایزگر کلاسیک که ورودی را به یک فضای خروجی محدود و گسسته (مانند یک برچسب مشخص) نگاشت می‌کنند، تولید متن فرایندی با ابعاد بالا، چندوجهی و تعاملی است~\cite{schneider2024explainable}. در این راستا، اشنایدر (۲۰۲۴)~\cite{schneider2024explainable} در تبیین چشم‌انداز هوش مصنوعی مولد تفسیرپذیر\enfootnote{generative explainable AI} (\lr{GenXAI}) استدلال می‌کند که به دلیل حجم عظیم اطلاعات در هر خروجی و فقدان یک حقیقت پایه، ارائه‌ی یک توضیح کامل و دقیق برای چرایی تولید یک توالی متنی، عملاً از نظر محاسباتی غیرقابل‌حل\enfootnote{intractable} است. موران و آراگام (۲۰۲۶)~\cite{moran2026towardinterpretable} نیز تأکید می‌کنند که مدل‌های مولد بازنمایی‌های خود را به‌شدت مبهم می‌سازند که تحلیل آن‌ها را به‌مراتب دشوارتر می‌کند. در این فضای پیوسته و درهم‌تنیده، حتی تغییرات جزئی در ورودی می‌تواند به خروجی‌های کاملاً متفاوت و غیرقابل‌پیش‌بینی منجر شود.

از منظر توپولوژیک و نظریه‌ی یادگیری آماری، این تمایز بنیادین میان مدل‌های متمایزگر و مولد را می‌توان در ساختار فضاهای هدف تبیین کرد. در یک مدل متمایزگر کلاسیک، تابع تصمیم $g: \mathcal{X} \to \Delta^{C-1}$ ورودی را به سیمپلکس\enfootnote{simplex} احتمالاتی روی مجموعه‌ای کوچک و متناهی از $C$ کلاس متمایز نگاشت می‌کند (یعنی $\Delta^{C-1} = \{ \mathbf{p} \in \mathbb{R}^C \mid \sum_{c=1}^C p_c = 1, p_c \ge 0 \}$)؛ این فضایی محدب با ابعاد بسیار پایین و ایستا است که در آن مفهوم «تصمیم» با شاخص بیشینه‌ی احتمال $\arg\max_c p_c$ هم‌ارز است. در مقابل، در مدل‌های زبانی مولد، فضای خروجی $\mathcal{Y} = \mathcal{V}^*$ یک فضای ترکیبیاتی با دامنه‌ی نامتناهی و ساختاری به‌شدت غیرخطی است؛ به‌گونه‌ای که حتی با محدود کردن طول توالی به $T_{\text{\lr{out}}}$ توکن، کاردینالیتی فضای خروجی برابر با $|\mathcal{V}|^{T_{\text{\lr{out}}}}$ است که به‌صورت نمایی با طول متن منفجر می‌شود. 

افزون بر این، نگاشت بازنمایی درونی ترنسفورمر $\phi_\theta: \mathcal{V}^* \to \mathcal{H}$ مفاهیم معنایی را در فضاهای پنهان چندبعدی به‌گونه‌ای درهم‌می‌آمیزد که تغییرات در فضای گسسته‌ی ورودی، رفتار پیوسته یا لیپشیتسی\enfootnote{Lipschitz continuous} از خود نشان نمی‌دهند؛ یعنی برای دو پرامپت با فاصله‌ی همینگ\enfootnote{Hamming distance} یک توکن ($d_H(\mathbf{x}, \mathbf{x}') = 1$)، توزیع‌های شرطی خروجی $P_\theta(\cdot \mid \mathbf{x})$ و $P_\theta(\cdot \mid \mathbf{x}')$ می‌توانند دارای واگرایی کولبک-لایبلر\enfootnote{Kullback-Leibler divergence} فوق‌العاده بزرگی باشند~\cite{zou2023universal}:
\begin{equation}\label{eq:chap1_kl_divergence}
    D_{\text{\lr{KL}}}\big(P_\theta(\cdot \mid \mathbf{x}) \parallel P_\theta(\cdot \mid \mathbf{x}')\big) \gg 0
\end{equation}
این عدم‌پایداری هندسی و فقدان یک تابع هدف با برچسب حقیقت پایه‌ی صریح ($\mathbf{y}^*$)، مسئله‌ی انتساب و تبیین رفتار مدل را از نظر محاسباتی به یک چالش دشوار ترکیبیاتی بدل می‌سازد.

از سوی دیگر، یکی از چالش‌های بنیادین در ادبیات فعلی پردازش زبان طبیعی، ابهام معنایی موجود پیرامون خود مفهوم انتساب\enfootnote{attribution} است. در ارزیابی مدل‌های زبانی مولد، رویکردهایی نظیر چارچوب «انتساب‌پذیر به منابع مشخص»\enfootnote{Attributable to Identified Sources} (\lr{AIS}) که توسط راشکین و همکاران (۲۰۲۳)~\cite{rashkin2023measuring} معرفی شد، مسئله‌ی انتساب را به معنای اعتبارسنجی خارجی و وفاداری به واقعیت تعریف می‌کنند؛ یعنی بررسی اینکه آیا گزاره‌ی تولیدشده در یک سند خارجی ریشه دارد یا خیر. اگرچه این رویکرد برای سنجش صحت خروجی و کاهش توهم اهمیت دارد، اما از تبیین سازوکار تصمیم‌گیری مدل در قبال ورودیِ ارائه‌شده بازمی‌ماند. خلط این دو مفهوم سبب شده است تا ارزیابی درستی خروجی، به‌اشتباه با تفسیر رفتار مدل هم‌ارز تلقی گردد. 

در نقطه‌ی مقابل، این پایان‌نامه انتساب را در قالب ارزیابی تجربی اهمیت ویژگی‌ها و ردیابی تأثیر اجزای پرامپت بر خروجی تعریف می‌کند. در این چارچوب، بخش‌ها یا جملاتی از پرامپت ورودی که نقشی در تولید خروجی ایفا می‌کنند را \textbf{«پیش‌آیند»} می‌نامیم و میزان وابستگی خروجی به حضور این گزاره‌ها را نمایانگر \textbf{«تأثیر»} آن‌ها بر رفتار مدل در نظر می‌گیریم. بر این اساس، مسئله‌ی انتساب بر سنجش این نکته متمرکز است که کدام پیش‌آیندهای موجود در پرامپت کاربر، بیشترین تأثیر را در تولید پاسخ هدف داشته‌اند. از این منظر، برخلاف چارچوب‌های اعتبارسنجی خارجی که در پی ارزیابی صدق گزاره نسبت به یک پایگاه دانش خارجی $\mathbb{I}(\mathbf{y} \models \mathcal{K}_{\text{\lr{ext}}})$ هستند، انتساب در این پژوهش بر پایه‌ی ارزیابی حساسیت رفتار خروجی مدل نسبت به حضور زیرمجموعه‌ای از جملات پرامپت ($\mathbf{x}_S$) در مقایسه با غیاب آن‌ها ($\mathbf{x} \setminus \mathbf{x}_S$) فرمول‌بندی می‌شود:
\begin{equation}\label{eq:chap1_causal_interv}
    \Delta_{\text{\lr{attr}}} = P_\theta(\mathbf{y}_T \mid \mathbf{x}_S) - P_\theta(\mathbf{y}_T \mid \mathbf{x} \setminus \mathbf{x}_S)
\end{equation}
این صورت‌بندی، میزان حساسیت و وابستگی رفتار مدل به گزاره‌های ورودی را مستقیماً ارزیابی می‌کند و کاملاً مستقل از صدق یا کذب گزاره در دنیای خارج است.

با این حال، پیش از ارائه‌ی صورت‌بندی ریاضی این مسئله، تعیین واحد تحلیل اهمیتی بنیادین دارد. روش‌های سنتی انتساب ویژگی در پردازش زبان طبیعی غالباً بر روی توکن‌های گسسته عمل می‌کنند؛ اما در مدل‌های زبانی خودهمبسته، توکن‌ها قطعات خُرد نحوی هستند که به‌تنهایی فاقد استقلال معنایی‌اند. حذف یا دستکاری توکن‌های منفرد، ساختار دستوری پرامپت را متلاشی کرده و داده‌هایی خارج از توزیع\enfootnote{out-of-distribution} (\lr{OOD}) پدید می‌آورد؛ وضعیتی که در آن واکنش مدل صرفاً بازتابی از تخریب گرامر است، نه نشانه‌ای از وابستگی محتوایی~\cite{madsen2022evaluating}. علاوه بر این، زبان طبیعی سرشار از افزونگی واژگانی و مترادف‌ها است؛ بنابراین، نوسان احتمال یک توکن خاص هرگز به معنای حذف بار معنایی یک مفهوم از خروجی نیست. 

برای رفع این چالش ساختاری، این پژوهش با اتکا به مبانی فلسفه‌ی زبان و علوم شناختی، واحد تحلیل را به سطح جملات ارتقا می‌دهد. بر مبنای «اصل بافتار»\enfootnote{context principle} گوتلوب فرگه (۱۸۸۴)~\cite{frege1991foundations}، واژگان منفرد در انزوا فاقد معنای مستقل هستند و محتوای حقیقی اندیشه تنها در افق یک ساختار گزاره‌ایِ\enfootnote{propositional} کامل شکل می‌گیرد. در امتداد همین بینش، والتر کینتش (۱۹۹۸)~\cite{kintsch1998comprehension} در مدل شناختیِ «ساخت‌وساز-ادغام» نشان می‌دهد که ذهن انسان متن را نه به‌عنوان توالی خرد کلمات، بلکه در قالب یک «پایگاه متن» شبکه‌بندی‌شده از گزاره‌ها پردازش می‌کند که بیانگر ایده‌ها و افکار کامل هستند. با انتخاب جمله به‌عنوان مفهوم و واحد عملیاتی تحلیل، تغییرات اعمال‌شده در ورودی هویت دستوری و معنایی خود را حفظ کرده و پیش‌آیندهای استخراج‌شده به‌صورت گزاره‌هایی معنادار، شفاف و انسان‌خوانا برای تحلیل‌گر تبیین می‌شوند. این مبنای نظری، پایه‌ریزی صوری فضای گزاره‌ها و عملگر تقطیع را در ادامه‌ی این بخش ایجاب می‌نماید.


\subsection{صورت‌بندی ریاضی و تعاریف پایه}\label{sec:1-1-1}

به‌منظور ارائه‌ی یک چارچوب دقیق و ساختاریافته بر پایه‌ی تحلیل در سطح گزاره‌ها، ابتدا فضاهای ورودی، خروجی و روابط تحلیلی مسئله را به‌صورت صوری پایه‌ریزی می‌کنیم.

\begin{definition}[فضای گزاره‌ها و عملگر تقطیع نحوی]\label{def:chap1_proposition_space}
فرض کنید $\mathcal{S} \subset \mathcal{V}^*$ فضای تمامی جملات (گزاره‌های زبانی با ساختار نحوی معتبر) باشد؛ به‌گونه‌ای که هر جمله خود از دنباله‌ای متوالی از توکن‌های الفبای $\mathcal{V}$ تشکیل شده است. عملگر تقطیع نحوی $\Pi: \mathcal{V}^* \to \bigcup_{k=1}^\infty \mathcal{S}^k$ تابعی خوش‌تعریف است که دنباله‌ی پیوسته‌ی توکن‌های یک متن را به دنباله‌ای مرتب و متناهی از جملات افراز می‌کند. بدین‌ترتیب، با در نظر گرفتن توکن‌های ورودی $(u_1, u_2, \ldots, u_{T_{\text{\lr{in}}}})$ و توکن‌های خروجی $(v_1, v_2, \ldots, v_{T_{\text{\lr{out}}}})$:
\begin{enumerate}
    \item پرامپت خام ورودی کاربر $\mathbf{x} = (u_1, u_2, \ldots, u_{T_{\text{\lr{in}}}}) \in \mathcal{V}^{T_{\text{\lr{in}}}}$ به دنباله‌ای از $N$ جمله نگاشت می‌شود:
    \begin{equation}\label{eq:chap1_prompt_decomp}
        \mathbf{x} = (x_1, x_2, \ldots, x_N) \in \mathcal{S}^N
    \end{equation}
    که در آن هر جمله‌ی ورودی $x_i \in \mathcal{S}$ خود از یک زیردنباله‌ی پیوسته از توکن‌های ورودی به فرم $x_i = (u_{\tau_{i-1}+1}, \ldots, u_{\tau_i})$ با شاخص‌های مرزی $0 = \tau_0 < \tau_1 < \cdots < \tau_N = T_{\text{\lr{in}}}$ متشکل است.
    \item پاسخ خام تولیدشده توسط مدل $\mathbf{y} = (v_1, v_2, \ldots, v_{T_{\text{\lr{out}}}}) \in \mathcal{V}^{T_{\text{\lr{out}}}}$ به دنباله‌ای از $M$ جمله نگاشت می‌شود:
    \begin{equation}\label{eq:chap1_response_decomp}
        \mathbf{y} = (y_1, y_2, \ldots, y_M) \in \mathcal{S}^M
    \end{equation}
    که در آن هر جمله‌ی خروجی $y_j \in \mathcal{S}$ خود از یک زیردنباله‌ی پیوسته از توکن‌های خروجی به فرم $y_j = (v_{\sigma_{j-1}+1}, \ldots, v_{\sigma_j})$ با شاخص‌های مرزی $0 = \sigma_0 < \sigma_1 < \cdots < \sigma_M = T_{\text{\lr{out}}}$ تشکیل یافته است.
\end{enumerate}
\end{definition}

\begin{definition}[زیرمجموعه‌ی هدف خروجی]\label{def:chap1_target_subset}
زیرمجموعه‌ی هدف خروجی که تحلیل‌گر قصد شناسایی پیش‌آیندهای اثرگذار بر آن را دارد، زیرمجموعه‌ای از شاخص‌های جملات خروجی $T \subseteq \{1, 2, \ldots, M\}$ با کاردینالیتی $|T| = M_T \le M$ است و به صورت توالی جملات هدف $\mathbf{y}_T$ تعریف می‌گردد:
\begin{equation}\label{eq:chap1_target_decomp}
    \mathbf{y}_T = {(y_j)}_{j \in T} \subseteq \mathbf{y}
\end{equation}
در حالت خاص که کاربر قصد تبیین منشأ یک گزاره‌ی منفرد را دارد، $T = \{t\}$ و $\mathbf{y}_T = y_t$ خواهد بود.
\end{definition}

\begin{definition}[زیرمجموعه‌ی پیش‌آیند ورودی]\label{def:chap1_mask_space}
فضای فرضیات انتساب، ابرمکعب دودویی یا مشبکه‌ی بولی\enfootnote{Boolean lattice} $\mathcal{M} = {\{0, 1\}}^N$ با بعد $N$ است. هر بردار ماسک دودویی $\mathbf{m} = {(m_1, m_2, \ldots, m_N)}^\top \in \mathcal{M}$ زیرمجموعه‌ای از جملات پرامپت ورودی را به‌عنوان پیش‌آیندهای گزینش‌شده مشخص می‌کند که با استفاده از عملگر ضرب هادامارد ($\odot$) (ضرب درایه‌به‌درایه) استخراج می‌گردد:
\begin{equation}\label{eq:chap1_subset_mask}
    \mathbf{x}_S = \mathbf{x} \odot \mathbf{m} = \{ x_i \in \mathbf{x} \mid m_i = 1, \, i \in \{1, \ldots, N\} \} \subseteq \mathbf{x}
\end{equation}
به‌طور متناظر، زیرمجموعه‌ی متمم یا دورریز به‌صورت زیر تعریف می‌شود:
\begin{equation}\label{eq:chap1_subset_discard}
    \mathbf{x}_{\text{\lr{discard}}} = \mathbf{x} \odot (\mathbf{1} - \mathbf{m}) = \{ x_i \in \mathbf{x} \mid m_i = 0, \, i \in \{1, \ldots, N\} \} = \mathbf{x} \setminus \mathbf{x}_S
\end{equation}
که در آن $\mathbf{1} = {(1, 1, \ldots, 1)}^\top \in {\{1\}}^N$ بردار یکه است.
\end{definition}

بر پایه‌ی تعاریف ارائه‌شده، صورت‌بندی دقیق مسئله‌ی بهینه‌سازی بنیادین این پژوهش در قالب مسئله‌ی~\ref{prob:chap1_causal_attribution} بیان می‌شود:

\begin{problem}[انتساب در سطح جمله برای مدل‌های زبانی جعبه‌سیاه]\label{prob:chap1_causal_attribution}
فرض کنید یک مدل زبانی جعبه‌سیاه $f_\theta: \mathcal{S}^N \to \mathcal{S}^M$ تنها از طریق استعلام مرتبه‌صفر در دسترس باشد، و جفت ورودی-خروجی $(\mathbf{x}, \mathbf{y})$ و زیرمجموعه‌ی هدف خروجی $\mathbf{y}_T \subseteq \mathbf{y}$ داده شده باشند. هدف، یافتن بردار ماسک دودویی بهینه‌ی $\mathbf{m}^* \in {\{0, 1\}}^N$ است که زیرمجموعه‌ی پیش‌آیند متناظر $\mathbf{x}_S = \mathbf{x} \odot \mathbf{m}^*$ را به‌گونه‌ای استخراج کند که هم‌زمان دو اصل بنیادین زیر ارضا شوند:
\begin{enumerate}
    \item \textbf{اصل کفایت\enfootnote{sufficiency}~\cite{ribeiro2018anchors, waldchen2021computational}:} اطلاعات موجود در $\mathbf{x}_S$ به‌تنهایی برای هدایت مدل به سوی بازتولید گزاره‌ی هدف $\mathbf{y}_T$ کافی باشد؛ به این معنا که میزان انحراف یا واگرایی معنایی پاسخِ حاصل از پرامپت تقلیل‌یافته نسبت به هدف خروجی از یک حد آستانه تجاوز نکند:
    \begin{equation}\label{eq:chap1_suff_condition}
        \mathcal{D}\big(f_\theta(\mathbf{x} \odot \mathbf{m}^*), \mathbf{y}_T\big) \le \epsilon
    \end{equation}
    که در آن $\mathcal{D}: \mathcal{S}^* \times \mathcal{S}^* \to \mathbb{R}_{\ge 0}$ یک شبه‌فاصله‌ی معنایی یا سنجه‌ی واگرایی در فضای خروجی مدل و $\epsilon > 0$ حد آستانه‌ی خطای پذیرفتنی است.
    
    \item \textbf{اصل ضرورت و خلوتی\enfootnote{necessity and sparsity}~\cite{chen2018learning, yoon2019invase, ignatiev2019abduction}:} زیرمجموعه‌ی $\mathbf{x}_S$ کمینه‌ترین (فشرده‌ترین) مجموعه‌ی گزاره‌های ممکن با ویژگی کفایت باشد؛ بدین معنا که هیچ پیش‌آیند غیرضروری یا زائدی در آن گنجانده نشده باشد و با حذف هر یک از اعضای آن، شرط کفایت نقض گردد:
    \begin{equation}\label{eq:chap1_sparsity_condition}
        \min_{\mathbf{m} \in {\{0, 1\}}^N} \|\mathbf{m}\|_0 = \sum_{i=1}^N m_i
    \end{equation}
    که در آن $\|\cdot\|_0$ شبه‌نرم $L_0$ بر روی فضای $\mathbb{R}^N$ است که تعداد مؤلفه‌های غیرصفر بردار را شمارش می‌کند.
\end{enumerate}
بنابراین، صورت‌بندی بهینه‌سازی ترکیبیاتی مقید مسئله به فرم صریح زیر خلاصه می‌شود~\cite{chen2018learning, waldchen2021computational}:
\begin{equation}\label{eq:chap1_constrained_opt}
    \mathbf{m}^* = \arg\min_{\mathbf{m} \in {\{0, 1\}}^N} \|\mathbf{m}\|_0 \quad \text{به‌شرط آنکه} \quad \mathcal{D}\big(f_\theta(\mathbf{x} \odot \mathbf{m}), \mathbf{y}_T\big) \le \epsilon
\end{equation}
\end{problem}


\subsection{تحلیل دشواری محاسباتی و بن‌بست‌های نظری مسئله}\label{sec:1-1-2}

صورت‌بندی ریاضی ارائه‌شده در معادله‌ی~\ref{eq:chap1_constrained_opt} پرده از سه چالش ساختاری عمیق برمی‌دارد که مواجهه‌ی سرراست با این مسئله را در مدل‌های زبانی مسدود می‌سازد:
\begin{enumerate}
    \item \textbf{پیچیدگی ترکیبیاتی نمایی فضای جستجو:} کاردینالیتی مشبکه‌ی بولی $\mathcal{M} = {\{0, 1\}}^N$ برابر با $|\mathcal{M}| = 2^N$ است. بهینه‌سازی تابعی بر روی این ابرمکعب گسسته بدون دسترسی به اطلاعات تحلیلی گرادیان، در حالت کلی در رده‌ی مسائل $\mathbf{NP}\text{\lr{-hard}}$ قرار دارد~\cite{waldchen2021computational, ignatiev2019abduction}. 
    
    به بیان دقیق‌تر، در تحلیل نظری مفسرهای یادگیری ماشین، ولدشن و همکاران (۲۰۲۱)~\cite{waldchen2021computational} و ایگناتیف و همکاران (۲۰۱۹)~\cite{ignatiev2019abduction} اثبات کرده‌اند که تصمیم‌گیری پیرامون وجود یک زیرمجموعه‌ی کافی از متغیرها با اندازه‌ی محدود در شبکه‌های عصبی یک مسئله‌ی $\mathbf{NP}\text{\lr{-complete}}$ است، و کشف کمینه‌ترین زیرمجموعه از منظر کاردینالیتی یک مسئله‌ی بهینه‌سازی ترکیبیاتی $\mathbf{NP}\text{\lr{-hard}}$ به شمار می‌رود.
    
    به‌عنوان مثال، برای یک پرامپت با طول متعارف ۳۰ جمله، فضای حالات شامل $2^{30} \approx 1.07 \times 10^9$ زیرمجموعه‌ی بالقوه است که ارزیابی مستقیم آن‌ها از نظر محاسباتی ناممکن است.
    
    \item \textbf{انسداد اوراکل در دسترسی جعبه‌سیاه:} تابع $f_\theta$ به‌صورت تحلیلی در دسترس نبوده و گرادیان آن قابل‌محاسبه نیست؛ بنابراین روش‌های بهینه‌سازی پیوسته (نظیر ضرب‌کننده‌های لاگرانژ یا روش‌های گرادیانی) مستقیماً بر روی $\mathbf{x} \odot \mathbf{m}$ قابل‌اعمال نیستند. علاوه بر این، در سیستم‌های تجاری، هر بار استعلام از اوراکل مستلزم صرف هزینه، پهنای باند و زمان تأخیر شبکه در بستر \lr{API} است که پرس‌وجوهای مکرر و بهینه‌سازی‌های مرتبه‌صفر را با محدودیت‌های شدید محاسباتی مواجه می‌سازد~\cite{nesterov2017random}.
    
    \item \textbf{عدم‌دسترسی به توزیع احتمالات خروجی و گرادیان‌ها:} در سیستم‌های بسته، خروجی اوراکل صرفاً متن نهایی است و بردار احتمالات توکن‌ها بر روی واژگان ($\mathbf{p} \in \Delta^{|\mathcal{V}|-1}$) ارائه نمی‌شود. این امر مانع از ارزیابی تغییرات پیوسته‌ی احتمالات شرطی شده و سنجه‌های انحراف باید مستقیماً بر روی متن محقق‌شده اعمال شوند.
\end{enumerate}

با در نظر گرفتن این چالش‌ها، پیچیدگی ابعاد بالای زبان طبیعی، و عدم‌دسترسی روش‌های متداول تفسیرپذیری جعبه‌سفید\enfootnote{white-box} و جعبه‌خاکستری\enfootnote{gray-box} به اطلاعات درونی مدل‌های تجاری مسدود، پرسش بنیادین این پژوهش شکل می‌گیرد:

\vspace{0.25cm}
\academicbox{مسئله‌ی اصلی پژوهش}{%
چگونه می‌توان بدون اعمال محدودیت در ساختار مدل‌های زبانی بزرگ و بدون دسترسی به سیگنال‌های داخلی (نظیر فعال‌سازی‌ها یا گرادیان‌ها)، رفتار یک سیستم مولد جعبه‌سیاه را تحلیل کرد و میزان تأثیر بخش‌های مختلف پرامپت ورودی را به‌عنوان پیش‌آیندهای اثرگذار بر روی یک زیرمجموعه‌ی خاص از خروجی، به شکلی ساختاریافته و انسان‌خوانا مشخص نمود؟
}
\vspace{0.25cm}


\section{اهمیت موضوع}\label{sec:1-2}

با نفوذ روزافزون هوش مصنوعی مولد به سیستم‌های تصمیم‌ساز و حوزه‌های پرمخاطره نظیر تشخیص‌های پزشکی، مشاوره‌ی حقوقی و قضاوت خودکار، شفافیت از یک ویژگی مطلوب به یک الزام اساسی تبدیل شده است. در این حوزه‌ها، خروجی‌های غیرشفاف و دارای سوگیری‌های پنهان\enfootnote{hidden biases} می‌توانند پیامدهای جانی، مالی یا اجتماعی جبران‌ناپذیری به همراه داشته باشند. دوشی-ولز و کیم (۲۰۱۷)~\cite{doshivelez2017rigorous} در مطالعه‌ای بنیادین استدلال می‌کنند که نیاز به تفسیرپذیری مستقیماً از نقص در فرمول‌بندی مسئله ناشی می‌شود؛ به این معنا که معیارهای کمّی مرسوم (مانند دقت در پیش‌بینی) برای تضمین اهداف حیاتی چون ایمنی، عدم‌سوگیری و جلب اعتماد انسان‌ها کافی نیستند. در محیط‌های حساس، عدم‌شفافیت مانع از آن می‌شود که متخصصان بتوانند فرایند اتخاذ تصمیم مدل را با دانش تخصصی خود اعتبارسنجی کنند یا پاسخگویی آن را بسنجند. این پژوهش تأکید می‌کند که برای استقرار مسئولانه‌ی این سیستم‌ها، نیازمند ارزیابی‌های کاربردمحور هستیم؛ هدفی که مدل‌های زبانی جعبه‌سیاه به دلیل ماهیت غیرشفاف خود دستیابی به آن را دشوار می‌سازند. از این‌رو، در چنین سیستم‌هایی دانستن منشأ اتخاذ یک تصمیم، هم‌ارز با خود آن تصمیم اهمیت می‌یابد.

این ضرورت با توجه به کیفیت ظاهری خروجی در مدل‌های زبانی پیشرفته، ابعاد حساس‌تری به خود می‌گیرد. مدل‌های زبانی بزرگ قادرند متونی با فصاحت بسیار بالا تولید کنند که می‌تواند خطاهای استدلالی یا توهمات سیستم را پشت ظاهری متقاعدکننده پنهان سازد~\cite{turpin2023language}. این سطح از شیوایی کلام، ریسک اعتماد کورکورانه به خروجی‌های نادرست را افزایش می‌دهد~\cite{huang2026hallucinations}. زمانی که مدل پاسخی نادرست را با لحنی قاطع و ظاهری مستند ارائه می‌دهد، تنها راه برای کشف خطا و جلوگیری از گمراهی، دسترسی به سازوکارهای تفسیرپذیری به‌منظور شناسایی پیش‌آیندهای مؤثر در شکل‌گیری آن متن است.

از منظر مهندسی نرم‌افزار و سامانه‌های هوشمند نیز، انتساب دقیق پیش‌شرطی اساسی برای کنترل و ارتقای مداوم مدل‌هاست. بدون درک نحوه‌ی اثرگذاری بخش‌های مختلف پرامپت بر خروجی، عیب‌یابی مدل‌ها با دشواری جدی مواجه خواهد شد~\cite{shbita2026systematic}. هنگامی که یک مدل زبانی رفتار نامطلوبی از خود بروز می‌دهد یا سوگیری مخربی را منعکس می‌سازد، ابزارهای انتساب از کارآمدترین روش‌ها برای شناسایی ریشه‌ی خطاها به‌شمار می‌روند~\cite{ijcai2024detecting}. فراتر از عیب‌یابی، هدایت مدل در مسیر پاسخگویی مسئولانه مستلزم آن است که بتوان الگوهای اثرگذاری ورودی بر خروجی را به شکلی نظام‌مند پایش کرد~\cite{lochab2025energybased, sharma2024acritical}.

در نهایت، تحولات در تدوین مقررات و چارچوب‌های حاکمیت هوش مصنوعی، تضمین شفافیت را به یک الزام قانونی تبدیل کرده است~\cite{kaminski2019right, papadimitriou2023right}. قوانین نوظهور، توسعه‌دهندگان را موظف می‌سازند تا امکان تبیین رفتار سیستم‌های هوشمند را فراهم آورند. مجموع این عوامل نشان می‌دهد که توسعه‌ی روش‌های نوین برای ممیزی مدل‌های مولد در شرایط جعبه‌سیاه، پاسخی به یکی از نیازهای ضروری اکوسیستم فناوری اطلاعات محسوب می‌شود.


\section{پیشینه و چالش‌های فعلی}\label{sec:1-3}

در این بخش، یک دسته‌بندی کلان از رویکردهای تفسیرپذیری ارائه می‌شود~\cite{guidotti2018survey, danilevsky2020survey} تا جایگاه و ضرورت روش پیشنهادی این پژوهش در میان پارادایم‌های موجود به روشنی تبیین گردد. فرمول‌بندی‌های تکمیلی و نقد تفصیلی این متدولوژی‌ها در فصل~\ref{chap:chapter3_literature} تشریح خواهد شد.

\subsection{تفسیرپذیری درونی}\label{sec:1-3-1}

رویکرد تفسیرپذیری درونی\enfootnote{intrinsic} بر طراحی و آموزش مدل‌هایی استوار است که ساختار و فرایند تصمیم‌گیری آن‌ها ذاتاً برای انسان شفاف باشد~\cite{danilevsky2020survey}. اگرچه این رویکرد در مدل‌های کلاسیک موفق بوده است – برای نمونه، روش‌هایی نظیر رگرسیون خطی تأثیر هر ویژگی را مستقیماً با یک ضریب مشخص نشان می‌دهند و درخت‌های تصمیم مسیر استنتاج را در قالب شروط صریح آشکار می‌سازند – اما در مواجهه با داده‌های پیچیده‌ای چون متون، این پارادایم با «مصالحه‌ی عملکرد-شفافیت»\enfootnote{performance-interpretability trade-off}~\cite{doshivelez2017rigorous} روبه‌رو می‌شود؛ بدین معنا که افزایش ظرفیت مدل جهت دستیابی به عملکرد بالاتر، معمولاً به بهای کاهش درک‌پذیری مستقیم سازوکارهای آن تمام می‌شود.

تلاش‌ها برای حفظ شفافیت در مدل‌های عمیق، به توسعه‌ی رویکردهایی نظیر یادگیری بازنمایی‌های تفکیک‌شده\enfootnote{disentangled representations} (به‌منظور تفکیک عوامل معنایی مستقل در ابعاد پنهان) و مدل‌های گلوگاه مفهومی\enfootnote{concept bottleneck models} (\lr{CBM})~\cite{koh2020cbm} انجامید که مدل را ملزم به تولید مفاهیم میانی تفسیرپذیر پیش از ارائه‌ی پاسخ نهایی می‌کنند. با این حال، اعمال چنین قیدهای ساختاری با ماهیت توزیع‌شده و درهم‌تنیده‌ی زبان طبیعی ناسازگار بوده و ظرفیت بیان‌گری مدل را کاهش می‌دهد. علاوه بر این، رویکرد درونی مستلزم بازطراحی و آموزش مدل از پایه است و بنابراین برای مدل‌های زبانی عظیم و پیش‌آموزش‌دیده‌ی کنونی که به‌صورت منجمد عرضه می‌شوند، قابل‌اعمال نیست.

\subsection{تفسیرپذیری پساعملی}\label{sec:1-3-2}

با توجه به محدودیت‌های رویکرد درونی، پژوهش‌ها به سوی پارادایم تفسیرپذیری پساعملی\enfootnote{post-hoc} معطوف شد. در این رهیافت، مدل ابتدا به‌طور کامل آموزش می‌بیند و سپس روش‌های ثانویه برای تفسیر رفتار مدلِ پیش‌آموزش‌دیده به‌کار گرفته می‌شوند. این حوزه بر اساس میزان دسترسی تحلیل‌گر به اطلاعات درونی شبکه، به دو شاخه‌ی اصلی جعبه‌سفید و جعبه‌سیاه تقسیم می‌شود~\cite{guidotti2018survey, danilevsky2020survey}.


\subsubsection{رویکردهای جعبه‌سفید}\label{sec:1-3-2-1}

روش‌های جعبه‌سفید بر این پیش‌فرض استوارند که تحلیل‌گر به اجزای گراف محاسباتی مدل، از جمله معماری، وزن‌ها و فضاهای فعال‌سازی، دسترسی دارد. این روش‌ها طیف وسیعی از سازوکارها را دربر می‌گیرند. در پایه‌ای‌ترین سطح، با ردیابی گرادیان‌ها (مانند روش گرادیان‌های یکپارچه یا \lr{IG}~\cite{sundararajan2017axiomatic}) و الگوهای توجه، ارتباط ویژگی‌های ورودی با خروجی بررسی می‌شود. فراتر از گرادیان‌های مستقیم، کاوش در فضای پنهان شبکه با بهره‌گیری از کاوشگرهای خطی\enfootnote{linear probes} و روش‌هایی چون \lr{TCAV}~\cite{kim2018interpretability} امکان استخراج مفاهیم هندسی را در لایه‌های میانی فراهم می‌سازد.

در سطوح ساختاری‌تر، رویکردهای تفسیرپذیری مکانیکی\enfootnote{mechanistic} و ویرایش مدل قرار دارند که با به‌کارگیری ابزارهایی نظیر خودرمزگذارهای خلوت\enfootnote{sparse autoencoders}~\cite{elhage2021mathematical} و نگاشت حالت‌های پنهان به توکن‌های زبانی~\cite{ghandeharioun2024patchscopes} در پی بررسی مدارهای محاسباتی درون مدل هستند. علاوه بر این، دسترسی به پارامترهای شبکه به پژوهشگران اجازه می‌دهد تا با دستکاری وزن‌ها (نظیر روش \lr{ROME}~\cite{meng2022locating})، دانش ذخیره‌شده را به‌صورت موضعی تغییر دهند.

با وجود غنای تئوریک، رویکردهای جعبه‌سفید با چالش‌های مشخصی مواجه‌اند. از یک سو، پژوهش‌های متعددی نشان داده‌اند که پراکسی‌های درونی همواره بازتاب کاملی از فرایند تصمیم‌گیری نیستند؛ از جمله اثبات عدم‌وفاداری مکانیزم توجه به فرایند تصمیم‌گیری~\cite{jain2019attention} و پدیده‌ی توهم تفسیرپذیری در برخی مداخلات مکانیکی~\cite{friedman2023comparing} که بیانگر فاصله‌ی بالقوه میان هندسه‌ی داخلی و مکانیزم واقعی اثرگذاری بر خروجی است. از سوی دیگر، در کاربردهای عملیاتی، قدرتمندترین مدل‌های زبانی به‌صورت سرویس‌های مسدود تجاری و در پسِ رابط‌های \lr{API} ارائه می‌شوند که دسترسی به وزن‌ها یا مشتقات را غیرممکن می‌سازد. با این حال، در شرایطی که دسترسی به مدل شفاف فراهم باشد، روش‌های جعبه‌سفید به دلیل بهره‌مندی مستقیم از گرادیان‌ها و اطلاعات درونی از مزیتی ساختاری و آشکار برخوردارند. از همین‌رو، در چارچوب ارزیابی‌های این پایان‌نامه نیز این روش‌ها به‌عنوان خط‌پایه‌های مقایسه‌ای در محیط‌های کنترل‌شده گنجانده شده‌اند تا سنجش جامعی از عملکرد چارچوب پیشنهادی در کنار روش‌های متکی بر اطلاعات درونی صورت پذیرد.

\subsubsection{رویکردهای جعبه‌سیاه}\label{sec:1-3-2-2}

در پارادایم جعبه‌سیاه، مدل صرفاً به‌عنوان یک نگاشت ورودی-خروجی بدون دسترسی به پارامترهای درونی در نظر گرفته می‌شود~\cite{guidotti2018survey}. این رهیافت رویکردی واقع‌گرایانه برای مواجهه با مدل‌های مسدود است. با این حال، روش‌های موجود در این دامنه با چالش‌های ساختاری مواجه‌اند که می‌توان آن‌ها را در چهار جریان اصلی دسته‌بندی کرد:

\paragraph{۱. روش‌های خودتوضیحی و استخراج متنی.}
نخستین نکته در این بخش، تفاوت ماهوی میان مسئله‌ی انتساب با رویکردهایی نظیر زنجیره‌ی افکار\enfootnote{chain-of-thought} (\lr{CoT})~\cite{wei2022chain} است. در مسئله‌ی انتساب، هدف نه تولید پیش‌رونده‌ی مراحل حل مسئله، بلکه شناسایی پس‌روندگانه‌ی میزان تأثیرگذاری بخش‌های ورودی بر خروجیِ حاصل‌شده است. شهودی‌ترین شیوه برای تعامل با مدل جعبه‌سیاه، استفاده از مهندسی پرامپت\enfootnote{prompt engineering} به‌منظور درخواست ارزیابی از خود مدل است. با این حال، مطالعات نشان داده‌اند که این توضیحات به پدیده‌ی استدلال هدفمند\enfootnote{motivated reasoning} دچار می‌شوند~\cite{turpin2023language}؛ به‌طوری‌که مدل‌ها به‌جای ارائه‌ی تصویری وفادار از تصمیم خود، به تولید توجیهات قانع‌کننده اما غیرمنطبق با منطق واقعی روی می‌آورند.

از سوی دیگر، استفاده از روش‌های بهینه‌سازی خودکار پرامپت (نظیر \lr{APE}~\cite{zhou2023large} و ابزارهای کنترل سوگیری~\cite{clemmer2024precisedebias}) یا متدولوژی‌های مبتنی بر حملات متخاصم\enfootnote{adversarial attacks} (مانند \lr{GCG}~\cite{zou2023universal} و \lr{PRP}~\cite{mangaokar2024prp}) برای گذر از سوگیری مدل و دسترسی به افکار داخلی، نیازمند یک خروجی هدفِ مشخص و ازپیش‌تعریف‌شده است؛ درحالی‌که در مسئله‌ی انتساب، حقیقت پایه‌ی معینی برای بهینه‌سازی وجود ندارد تا بتوان از طریق آن پرامپت‌های بهینه برای استخراج توضیحات وفادارانه تولید کرد.

\paragraph{۲. انتساب ویژگی از طریق اختلال در ورودی.}
جریان دوم که مبنای بسیاری از روش‌های جعبه‌سیاه به‌شمار می‌رود، بر پایه‌ی اعمال تغییرات در ورودی استوار است. در چارچوب‌های کلاسیک (نظیر \lr{LIME}~\cite{ribeiro2016lime} و \lr{SHAP}~\cite{lundberg2017unified})، الگوریتم تلاش می‌کند با دستکاری یا حذف بخش‌هایی از ورودی و مشاهده‌ی تغییرات خروجی، میزان تأثیر هر بخش را تخمین بزند. با این حال، تعمیم این سازوکار به مدل‌های زبانی با چالش‌های زیر همراه است:
\begin{itemize}
    \item \textbf{وابستگی به توزیع احتمالات خروجی:} بسیاری از روش‌های اختلالی برای محاسبه‌ی تغییرات خروجی به احتمالات توکن‌ها وابسته هستند، در حالی که در بسیاری از رابط‌های تجاری تنها متن نهایی بازگردانده می‌شود.
    \vspace{-0.15cm}
    \item \textbf{فروپاشی نحوی در اختلال‌های سطح توکن:} حذف توکن‌های مجزا موجب تخریب گرامر پرامپت شده و ورودی‌هایی \lr{OOD} ایجاد می‌کند. در این شرایط، مشخص نیست که افت کیفیت خروجی ناشی از حذف یک بخش محتوایی است یا صرفاً واکنش به تخریب ساختار دستوری متن~\cite{madsen2022evaluating}.
    \vspace{-0.15cm}
    \item \textbf{ابهام معنایی و هم‌پوشانی واژگانی:} در تولید زبان طبیعی توسط \lr{LLM}ها، یک مفهوم را می‌توان با کلمات گوناگونی بیان کرد. از این‌رو، افت احتمال یک توکن منفرد لزوماً به معنای حذف آن مفهوم از خروجی نیست؛ این در حالی است که در دسته‌بندی، یک کلمه نماینده‌ی یک معنای مستقل است.
    \vspace{-0.15cm}
    \item \textbf{پیچیدگی محاسباتی استعلام‌ها:} یافتن ترکیب بهینه‌ی اجزای ورودی مستلزم بررسی ترکیبیاتی ($\mathcal{O}(2^N)$) زیرمجموعه‌ها است که تکرار آن از طریق استعلام‌های \lr{API} برای پرامپت‌های طولانی هزینه و زمان غیرعملی به همراه دارد.
\end{itemize}

\paragraph{۳. یادگیری مدل‌های جانشین مفسر.}
برای کاستن از هزینه‌ی استعلام‌های مکرر در زمان اجرا، استفاده از مدل‌های جانشین\enfootnote{surrogate} پیشنهاد شده است (نظیر روش‌های \lr{L2X}~\cite{chen2018learning} و \lr{SOInter}~\cite{seyyedsalehi2024sointer}). در این پارادایم، یک مدل ثانویه آموزش می‌بیند تا رفتار مدل هدف را تقریب بزند. با این حال، اغلب فرمول‌بندی‌های موجود برای فضاهای خروجی با ابعاد ثابت (نظیر دسته‌بندی) طراحی شده‌اند و تعمیم آن‌ها به فضای نامتناهی و گسسته‌ی متن آزاد با دشواری‌های ساختاری روبه‌روست.

\paragraph{۴. چالش‌های ارزیابی و فقدان حقیقت پایه.}
در نهایت، ارزیابی وفاداری روش‌های انتساب در مدل‌های جعبه‌سیاه با چالش فقدان حقیقت پایه مواجه است. تلاش‌ها برای استانداردسازی ارزیابی‌ها (مانند مجموعه‌داده‌ی محک \lr{ERASER}~\cite{deyoung2020eraser}) نشان داده‌اند که معیارهای فرسایشی گاه دچار سوگیری‌های ارزیابی یا «قانون گودهارت»\enfootnote{Goodhart's law} می‌شوند~\cite{ju2022logictraps, hsia2024goodhart}؛ به‌طوری‌که مفسرها به‌جای انتساب گزاره‌های واقعاً اثرگذار، کلماتی را استخراج می‌کنند که بیشترین شوک توزیعی (\lr{OOD}) را ایجاد نمایند.

مجموع این نارسایی‌ها ضرورت بازنگری در واحد تحلیل، شیوه‌ی مدل‌سازی و نحوه‌ی تقریب رفتار را در شرایط جعبه‌سیاه آشکار می‌سازد.


\section{رویکرد پیشنهادی}\label{sec:1-4}

برای غلبه بر بن‌بست‌های نظری و محاسباتی رویکردهای جعبه‌سیاه (نظیر فروپاشی توزیعی در اثر حذف توکن‌ها، ناتوانی در ردیابی تغییرات معنادار به دلیل افزونگی زبانی، نیاز مبرم به توزیع احتمالات و لوجیت‌های مدل هدف، و هزینه‌ی نمایی استعلام‌ها)، این پایان‌نامه یک تغییر پارادایم اساسی در واحد تحلیل و شیوه‌ی تقریب ارائه می‌دهد. رویکرد پیشنهادی با ترکیب علوم شناختی، فلسفه‌ی زبان و مدل‌سازی انرژی، تلاش می‌کند تا بر محدودیت‌های کاربردی و ماهوی شبکه‌های مولد فائق آید. در ادامه‌ی این بخش، یک معرفی کلان و ساختاریافته از این رویکرد ارائه می‌شود؛ جزئیات فنی و فرمول‌بندی دقیق این معماری در فصل~\ref{chap:chapter4_methodology} به‌تفصیل تشریح خواهد شد.

نخستین مؤلفه در این رویکرد، \textbf{تحلیل در سطح گزاره‌ها و جملات} به‌عنوان مفاهیم معنایی پایدار است. همان‌طور که پیش‌تر در بخش~\ref{sec:1-1} بر پایه‌ی مبانی فلسفی و شناختی تبیین شد، روش‌های اختلالی مرسوم با دستکاری توکن‌ها، ساختار نحوی پرامپت را متلاشی کرده و مدل را با ورودی‌های \lr{OOD} مواجه می‌سازند. علاوه بر این، در متون طبیعی به دلیل افزونگی واژگانی، افت احتمال یک توکن خاص لزوماً به معنای محو آن مفهوم از خروجی نیست. با استناد به اصل بافتار فرگه~\cite{frege1991foundations} و نظریه‌ی شناختی کینتش~\cite{kintsch1998comprehension}، ارتقای واحد تحلیل به جمله تضمین می‌کند که حذف یا انتخاب یک ویژگی در فرایند ارزیابی، ماهیت ساختاری و معنایی خود را حفظ کرده و مفسر صرفاً واکنش شبکه به فروپاشی گرامر را اندازه‌گیری نمی‌کند. افزون بر این، چنین رویکردی کیفیت و شفافیت انتساب را برای کاربر انسانی به‌مراتب ارتقا می‌بخشد؛ زیرا ذهن انسان نقش و تأثیر یک گزاره‌ی کامل را بسیار واضح‌تر و شهودی‌تر از توکن‌های پراکنده و منقطع درک می‌نماید. بدین‌ترتیب، جمله به‌عنوان مفهوم عملیاتی بنیادین برای شناسایی پیش‌آیندها تعریف می‌گردد.

دومین مؤلفه در این معماری، گذر از تقریب‌های محلی و \textbf{مدل‌سازی سراسری الگوهای انطباق‌پذیری} است. مدل‌های مفسر محلی (نظیر \lr{LIME}) صرفاً به همسایگی‌های بسیار محدود پیرامون نمونه می‌نگرند که این امر غالباً به بروز توهم تفسیرپذیری و ناتوانی در درک تعاملات گسترده‌تر متنی منجر می‌شود~\cite{friedman2024interpretability}. از این‌رو، این پژوهش از پارادایمِ «آموزش سراسری یک مفسر محلی»~\cite{chen2018learning} بهره می‌برد تا الگوهای کلان هم‌نشینی متن را کشف کند. برای تحقق این امر، به‌جای تلاش برای نفوذ به لایه‌های پنهان شبکه‌ی هدف یا استعلام مکرر لوجیت‌ها، از یک مدل مبتنی بر انرژی\enfootnote{energy-based model} (\lr{EBM}) مبتنی بر معماری ترنسفورمر به‌عنوان جانشین استفاده می‌شود.

انتخاب مدل‌های مبتنی بر انرژی در این متدولوژی، علاوه بر مرتفع ساختن نیاز به دسترسی به لوجیت‌ها، پاسخی ساختاری به چالش ارزیابی در زبان طبیعی است. برخلاف مدل‌های خودهمبسته‌ی استاندارد که متن را به‌صورت ترتیبی (توکن‌به‌توکن) تولید و ارزیابی می‌کنند و به نرمال‌سازی احتمالات در فضای نمایی خروجی‌ها وابسته‌اند، یک مدل انرژی به‌عنوان یک تابع انطباق‌پذیریِ نرمال‌نشده\enfootnote{unnormalized compatibility function} عمل می‌کند. این معماری می‌تواند کل توالی را به‌طور یکپارچه و کل‌نگر بررسی کرده و میزان انطباق‌پذیری ساختاری و معنایی میان پرامپت و پاسخ مشاهده‌شده را در قالب یک امتیاز اسکالر در یک چشم‌انداز انرژی\enfootnote{energy landscape} ارزیابی نماید. این مدل با آموزش بر روی جفت‌های اصیل پرامپت-پاسخ و مقایسه‌ی آن‌ها با نسخه‌های مخدوش‌شده، الگوهای ساختاری و معنایی هم‌نشینی میان ورودی و خروجی را که در متون تولیدشده توسط مدل هدف مشاهده می‌شوند، فرامی‌گیرد. مهم‌تر از همه، این فرایند تنها نیازمند متن نهایی تولیدشده توسط سیستم جعبه‌سیاه است و وابستگی به معماری داخلی، گرادیان‌ها یا لوجیت‌ها ندارد. در این فضا، مقادیر پایین‌تر انرژی نمایانگر انطباق‌پذیری بالاتر بر اساس الگوهای مشاهده‌شده در متون تولیدی مدل است، بدون آنکه فرضی پیرامون یادگیری خمینه‌ها یا پویایی‌های بازنمایی درونی مدل هدف اعمال شود.

در نهایت، با تلفیق این نوآوری‌ها، این پایان‌نامه چارچوب \textbf{«جانشین مبتنی بر انرژی برای تفسیر در سطح مفهوم»} یا \lr{ESCI}\enfootnote{Energy-Based Surrogate for Concept-Level Interpretation} را معرفی می‌نماید. در این ساختار، پس از آنکه مدل مبتنی بر انرژی الگوهای انطباق‌پذیری میان گزاره‌های ورودی و خروجی را فراگرفت، سیگنال‌های ارزیابی آن برای آموزش یک شبکه‌ی مفسر اختصاصی به‌کار گرفته می‌شوند. در فاز عملیاتی (استنتاج)، با دریافت پرامپت ورودی ($\mathbf{x}$) و پاسخ تولیدشده ($\mathbf{y}$) از سوی مدل زبانی، کاربر یک زیرمجموعه‌ی هدف از خروجی ($\mathbf{y}_T \subseteq \mathbf{y}$) را مشخص می‌کند. سپس مفسر آموزش‌دیده با تخصیص وزن‌های اهمیت، زیرمجموعه‌ای از جملات پرامپت ($\mathbf{x}_S \subseteq \mathbf{x}$) را استخراج می‌کند که بیشترین تأثیر را به‌عنوان پیش‌آیند در شکل‌گیری آن بخش از خروجی داشته‌اند. این چارچوب با کار بر روی متن خام و حذف وابستگی به فراخوانی‌های مکرر \lr{API} در زمان استنتاج، انتساب پساعملی در سطح جمله را برای مدل‌های تجاری در شرایط جعبه‌سیاه به شکلی کارآمد امکان‌پذیر می‌سازد.

بدین‌ترتیب، رویکرد ارائه‌شده در این پژوهش اساساً بر ارزیابی میزان انطباق‌پذیری ساختاری و معنایی میان گزاره‌ها در بستر متون تولیدشده استوار است و از این طریق، جملات پیش‌آیندی را که بیشترین همخوانی را در شکل‌گیری بخش هدف خروجی دارند، شناسایی می‌نماید. اگرچه این رهیافت تجربی بر پایه‌ی الگوهای متنیِ مشاهده‌پذیر عمل می‌کند و متمایز از آزمون‌های مداخله‌ای یا اثبات‌های صوری علیت است، با این وجود نگاشت‌های انطباق‌پذیریِ به‌دست‌آمده از مدل مبتنی بر انرژی، نشانگرهایی تجربی و عینی از نحوه‌ی اثرگذاری ورودی بر پاسخ فراهم می‌آورند که می‌توانند سرنخ‌هایی معتبر و همسو با خطوط زیربنایی تأثیرگذاری میان جملات در مدل‌های زبانی جعبه‌سیاه ارائه دهند.


\section{اهداف پژوهش}\label{sec:1-5}

هدف کلان این پایان‌نامه، طراحی و توسعه‌ی یک چارچوب انتساب پساعملی و مستقل از مدل\enfootnote{model-agnostic} در سطح مفهوم برای مدل‌های زبانی بزرگ با دسترسی جعبه‌سیاه است؛ چارچوبی که بتواند جملات اثرگذار ورودی (پیش‌آیندها) را در شکل‌گیری بخش معینی از پاسخ تولیدشده، تنها بر مبنای تحلیل متن خام خروجی و بدون اتکا به وزن‌ها، گرادیان‌ها یا لوجیت‌ها شناسایی کند. با ارتقای واحد تحلیل به گزاره‌های کامل و به‌کارگیری مدل‌های مبتنی بر انرژی، این پژوهش سازوکاری منسجم برای سنجش وابستگی‌های گزاره‌ای در شرایط دسترسی محدود ارائه می‌دهد. به‌منظور سنجش قابلیت پیاده‌سازی، کیفیت انتساب و کارایی محاسباتی، چارچوب پیشنهادی در محیط‌های تجربی کنترل‌شده هم در برابر روش‌های خط‌پایه‌ی جعبه‌سیاه و هم در مقایسه با روش‌های جعبه‌سفید و مبتنی بر گرادیان (با در نظر داشتن مزیت آشکار ناشی از دسترسی آن‌ها به لایه‌های پنهان و مشتقات درونی شبکه) ارزیابی و اعتبارسنجی می‌شود.

برای تحقق این هدف کلان، اهداف فنی و ساختاری زیر دنبال می‌شوند:
\begin{enumerate}
    \item \textbf{طراحی تابع زیان آگاه به ساختار برای مدل انرژی:} آموزش شبکه‌ی انرژی با استفاده از یک تابع زیان دوگانه، به‌گونه‌ای که ارزیابی انطباق‌پذیری سراسری میان پرامپت و پاسخ را با ارزیابی وابستگی‌های محلی در سطح جملات تلفیق کند. این طراحی تضمین می‌کند که ارزیابی انطباق‌پذیری، بر پایه‌ی ارتباطات معنادار گزاره‌ای استوار شده و از اتکا به همبستگی‌های سطحی (نظیر ایست‌واژه‌ها) اجتناب گردد.
    
    \item \textbf{کاهش پیچیدگی محاسباتی در فاز استنتاج:} برخلاف روش‌های اختلالی سنتی که مستلزم ارسال استعلام‌های متعدد به رابط‌های برنامه‌نویسی هستند، در این معماری با آموزش یک شبکه‌ی مفسر اختصاصی جانشین بر مبنای سیگنال‌های مدل انرژی، بار محاسباتی انتساب در زمان اجرا به یک ارزیابی مستقیم و زمان‌ثابت تقلیل می‌یابد و وابستگی به استعلام‌های مکرر از مدل هدف به‌طور کامل برطرف می‌شود.
\end{enumerate}


\section{ساختار پایان‌نامه}\label{sec:1-6}

ساختار این پایان‌نامه در شش فصل تدوین شده است تا مسیری منطقی از بیان مفاهیم پایه تا ارزیابی روش پیشنهادی را طی کند. پس از این مقدمه، در \textbf{فصل دوم}، \nameref{chap:chapter2_preliminaries} و پایه‌های نظری مورد نیاز شامل معماری مدل‌های زبانی بزرگ، بنیان‌های ریاضی مدل‌های مبتنی بر انرژی و روش‌های یادگیری تقابلی تشریح می‌شوند. سپس در \textbf{فصل سوم}، \nameref{chap:chapter3_literature} و ادبیات موضوع به‌تفصیل واکاوی می‌گردد و رویکردهای موجود در تفسیرپذیری همراه با مزایا و چالش‌های نظری و عملی آن‌ها در مواجهه با سیستم‌های مدرن بررسی می‌شوند. در \textbf{فصل چهارم}، \nameref{chap:chapter4_methodology} این پژوهش (\lr{ESCI}) به‌صورت جامع معرفی شده و ساختار شبکه‌های انرژی و مفسر، به همراه الگوریتم‌های آموزش آن‌ها فرمول‌بندی می‌گردد. \textbf{فصل پنجم} به \nameref{chap:chapter5_experiments} اختصاص دارد؛ در این فصل، چارچوب پیشنهادی در محیط‌های تجربی کنترل‌شده و در مقایسه با خط‌پایه‌های جعبه‌سیاه و نیز خط‌پایه‌های جعبه‌سفید و مبتنی بر گرادیان، از منظر کیفیت انتساب، وفاداری تجربی و کارایی محاسباتی مورد سنجش قرار می‌گیرد. در نهایت، در \textbf{فصل ششم}، \nameref{chap:chapter6_conclusions}، جمع‌بندی دستاوردهای پژوهش ارائه شده و ضمن بررسی ابعاد کاربردی و محدودیت‌های موجود، مسیرهایی برای تحقیقات آتی پیشنهاد می‌گردد.
